\documentclass[10pt,a4paper]{amsart}
\usepackage{amsmath,amssymb,amsthm,graphicx}
% Retained minimal preamble for the Stacks Project chapter more-algebra.tex (More on Algebra).
% The chapter file pinned-original-source.tex carries only \input{preamble} on its first line, so no
% preamble text lives inside the chapter itself.  The Stacks Project uses one single shared preamble for
% every chapter; the snapshot stages it identically as sources/stacks-chow/preamble.tex and
% sources/stacks-etale-cohomology/preamble.tex.  The block below is copied mechanically, line for line,
% from sources/stacks-chow/preamble.tex lines 14-18 (the Xy-pic setup) and 154-217 (the native theorem
% declarations together with the mathematical macros).  Nothing was added, removed or redefined; the
% wrapper class and the amsmath/amssymb/amsthm/graphicx base packages are supplied by the assembler.
% Chapter: more-algebra.tex -- More on Algebra.
% For commutative diagrams we use Xy-pic
\usepackage[all]{xy}

% We use 2cell for 2-commutative diagrams.
\xyoption{2cell}

% Theorem environments.
%
\theoremstyle{plain}
\newtheorem{theorem}[subsection]{Theorem}
\newtheorem{proposition}[subsection]{Proposition}
\newtheorem{lemma}[subsection]{Lemma}

\theoremstyle{definition}
\newtheorem{definition}[subsection]{Definition}
\newtheorem{example}[subsection]{Example}
\newtheorem{exercise}[subsection]{Exercise}
\newtheorem{situation}[subsection]{Situation}

\theoremstyle{remark}
\newtheorem{remark}[subsection]{Remark}
\newtheorem{remarks}[subsection]{Remarks}

\numberwithin{equation}{subsection}

% Macros
%
\def\lim{\mathop{\mathrm{lim}}\nolimits}
\def\colim{\mathop{\mathrm{colim}}\nolimits}
\def\Spec{\mathop{\mathrm{Spec}}}
\def\Hom{\mathop{\mathrm{Hom}}\nolimits}
\def\Ext{\mathop{\mathrm{Ext}}\nolimits}
\def\SheafHom{\mathop{\mathcal{H}\!\mathit{om}}\nolimits}
\def\SheafExt{\mathop{\mathcal{E}\!\mathit{xt}}\nolimits}
\def\Sch{\mathit{Sch}}
\def\Mor{\mathop{\mathrm{Mor}}\nolimits}
\def\Ob{\mathop{\mathrm{Ob}}\nolimits}
\def\Sh{\mathop{\mathit{Sh}}\nolimits}
\def\NL{\mathop{N\!L}\nolimits}
\def\CH{\mathop{\mathrm{CH}}\nolimits}
\def\proetale{{pro\text{-}\acute{e}tale}}
\def\etale{{\acute{e}tale}}
\def\QCoh{\mathit{QCoh}}
\def\Ker{\mathop{\mathrm{Ker}}}
\def\Im{\mathop{\mathrm{Im}}}
\def\Coker{\mathop{\mathrm{Coker}}}
\def\Coim{\mathop{\mathrm{Coim}}}

% Boxtimes
%
\DeclareMathSymbol{\boxtimes}{\mathbin}{AMSa}{"02}

%
% Macros for moduli stacks/spaces
%
\def\QCohstack{\mathcal{QC}\!\mathit{oh}}
\def\Cohstack{\mathcal{C}\!\mathit{oh}}
\def\Spacesstack{\mathcal{S}\!\mathit{paces}}
\def\Quotfunctor{\mathrm{Quot}}
\def\Hilbfunctor{\mathrm{Hilb}}
\def\Curvesstack{\mathcal{C}\!\mathit{urves}}
\def\Polarizedstack{\mathcal{P}\!\mathit{olarized}}
\def\Complexesstack{\mathcal{C}\!\mathit{omplexes}}
% \Pic is the operator that assigns to X its picard group, usage \Pic(X)
% \Picardstack_{X/B} denotes the Picard stack of X over B
% \Picardfunctor_{X/B} denotes the Picard functor of X over B
\def\Pic{\mathop{\mathrm{Pic}}\nolimits}
\def\Picardstack{\mathcal{P}\!\mathit{ic}}
\def\Picardfunctor{\mathrm{Pic}}
\def\Deformationcategory{\mathcal{D}\!\mathit{ef}}
\usepackage{cleveref}
\crefname{originallemma}{original lemma}{original lemmas}
\Crefname{originallemma}{Original lemma}{Original lemmas}
\makeatletter
\newlabel{more-algebra-lemma-flat-unramified}{{0AGX}{source}{}{sourcepointer-1}{}}
\newlabel{more-algebra-lemma-flat-unramified@cref}{{[originallemma][0][]0AGX}{[1][0][]source}{}{}{}}
\makeatother
\crefname{originallemma}{original lemma}{original lemmas}
\Crefname{originallemma}{Original lemma}{Original lemmas}
\makeatletter
\newlabel{algebra-lemma-finite-after-completion}{{0394}{source}{}{sourcepointer-2}{}}
\newlabel{algebra-lemma-finite-after-completion@cref}{{[originallemma][0][]0394}{[1][0][]source}{}{}{}}
\makeatother
\crefname{originallemma}{original lemma}{original lemmas}
\Crefname{originallemma}{Original lemma}{Original lemmas}
\makeatletter
\newlabel{lemma-flat-after-completion}{{0AGW}{source}{}{sourcepointer-3}{}}
\newlabel{lemma-flat-after-completion@cref}{{[originallemma][0][]0AGW}{[1][0][]source}{}{}{}}
\makeatother
\crefname{originallemma}{original lemma}{original lemmas}
\Crefname{originallemma}{Original lemma}{Original lemmas}
\makeatletter
\newlabel{algebra-lemma-finite-flat-local}{{00NZ}{source}{}{sourcepointer-4}{}}
\newlabel{algebra-lemma-finite-flat-local@cref}{{[originallemma][0][]00NZ}{[1][0][]source}{}{}{}}
\makeatother
\crefname{originallemma}{original lemma}{original lemmas}
\Crefname{originallemma}{Original lemma}{Original lemmas}
\makeatletter
\newlabel{algebra-lemma-hathat-finitely-generated}{{05GG}{source}{}{sourcepointer-5}{}}
\newlabel{algebra-lemma-hathat-finitely-generated@cref}{{[originallemma][0][]05GG}{[1][0][]source}{}{}{}}
\makeatother
\title{A flat local homomorphism that is unramified on residue fields induces an isomorphism of completions}
\author{The Stacks Project Authors}
\date{stacks-project:more-algebra.tex:local-snapshot-20261012; GNU FDL 1.2 or later}
\begin{document}
\maketitle
\setcounter{subsection}{1}
\noindent\textit{Editorial note.} All delivered mathematical text is the author's own, copied from the pinned Stacks Project chapter file more-algebra.tex; no statement or proof step was rewritten, reordered or removed.\par\medskip
\noindent\textit{Editorial note.} The chapter file pinned-original-source.tex carries only \textbackslash{}input\{preamble\} on its first line, so no preamble text lives inside the chapter. The retained file sources/stacks-more-algebra/retained-preamble.tex is a mechanical, line-for-line copy of the project's single shared preamble as staged in the snapshot (Xy-pic setup plus theorem declarations and mathematical macros). Nothing was added, removed or redefined, and the excerpt serves only as preamble, never as a span coordinate.\par\medskip
\noindent\textit{Editorial note.} Both counted spans start at the \textbackslash{}begin line of their environment, so the statement and the author's proof are delivered whole; surrounding prose of the section is deliberately outside the counted spans and is carried as standard prerequisites instead.\par\medskip
\noindent\textit{Editorial note.} The citation fields (source\_id, source\_version, authors, source\_url, license.name, source\_title) are identical to those of every other card this lane has produced from this chapter, so the batch selector treats them as one source.\par\medskip
\noindent\textit{Allowed standard prerequisites (editorial).}\begin{itemize}
\item Flatness of a local homomorphism and the fibre criterion for flatness. used at proof lines 10772-10776. Supplies the flatness hypothesis and the way it is transported to the completion.
\item Finiteness and flatness over a Noetherian local ring imply freeness. the author cites Algebra, Lemma algebra-lemma-finite-flat-local at proof lines 10780-10782. Turns flatness into freeness of the completed map.
\item Completions of Noetherian local rings and the identification of the residue fields of the source and target. used at proof lines 10782-10787 together with the hypothesis on the residue fields. Provides the rank computation that shows the free module has rank one.
\end{itemize}\medskip
\ifdefined\hypertarget\hypertarget{sourcepointer-1}{}\fi
\noindent\textit{Source reference.} Original lemma 0AGX: focused-seeds/AGP/sources/stacks-snapshot-20261011/tags.txt, source lines 13210--13210. The lemma's own canonical tag, quoted so that readers can resolve the label inside the delivered text; it points at the counted statement itself and carries no proof content.\par\medskip
\ifdefined\hypertarget\hypertarget{sourcepointer-2}{}\fi
\noindent\textit{Source reference.} Original lemma 0394: focused-seeds/AGP/sources/stacks-snapshot-20261011/tags.txt, source lines 4259--4259. Navigation pointer to a statement the counted proof refers to but does not contain; nothing load-bearing is delegated to it.\par\medskip
\ifdefined\hypertarget\hypertarget{sourcepointer-3}{}\fi
\noindent\textit{Source reference.} Original lemma 0AGW: focused-seeds/AGP/sources/stacks-snapshot-20261011/tags.txt, source lines 13209--13209. Navigation pointer to a statement the counted proof refers to but does not contain; nothing load-bearing is delegated to it.\par\medskip
\ifdefined\hypertarget\hypertarget{sourcepointer-4}{}\fi
\noindent\textit{Source reference.} Original lemma 00NZ: focused-seeds/AGP/sources/stacks-snapshot-20261011/tags.txt, source lines 877--877. Navigation pointer to a statement the counted proof refers to but does not contain; nothing load-bearing is delegated to it.\par\medskip
\ifdefined\hypertarget\hypertarget{sourcepointer-5}{}\fi
\noindent\textit{Source reference.} Original lemma 05GG: focused-seeds/AGP/sources/stacks-snapshot-20261011/tags.txt, source lines 7060--7060. Navigation pointer to a statement the counted proof refers to but does not contain; nothing load-bearing is delegated to it.\par\medskip
\begin{lemma}
\label{lemma-flat-unramified}
Let $A \to B$ be a flat local homomorphism of Noetherian local rings
such that $\mathfrak m_A B = \mathfrak m_B$ and
$\kappa(\mathfrak m_A) = \kappa(\mathfrak m_B)$.
Then $A \to B$ induces an isomorphism $A^\wedge \to B^\wedge$
of completions.
\end{lemma}

\begin{proof}
By Algebra, Lemma \ref{algebra-lemma-finite-after-completion} we see that
$B^\wedge$ is the $\mathfrak m_A$-adic
completion of $B$ and that $A^\wedge \to B^\wedge$ is finite.
Since $A \to B$ is flat we have $\text{Tor}_1^A(B, \kappa(\mathfrak m_A)) = 0$.
Hence we see that $B^\wedge$ is flat over $A^\wedge$ by
Lemma \ref{lemma-flat-after-completion}.
Thus $B^\wedge$ is a free $A^\wedge$-module by
Algebra, Lemma \ref{algebra-lemma-finite-flat-local}.
Since $A^\wedge \to B^\wedge$ induces an isomorphism
$\kappa(\mathfrak m_A) = A^\wedge/\mathfrak m_A A^\wedge \to
B^\wedge/\mathfrak m_A B^\wedge = B^\wedge/\mathfrak m_B B^\wedge =
\kappa(\mathfrak m_B)$ by our assumptions
(and Algebra, Lemma \ref{algebra-lemma-hathat-finitely-generated}),
we see that $B^\wedge$ is free of rank $1$. Thus $A^\wedge \to B^\wedge$
is an isomorphism.
\end{proof}

\end{document}
